\section{PIの中心幾何から得られる有限生成子}
\label{sec:notes-pi-residue-generator}

有限位数のSklyaninブロックでは、本文の外部積生成定理とは別に、
中心の零sliceのAzumaya幾何から全MF圏のpropernessと有限生成子を直接得られる。
この証明は \(\gcd(|\sigma_i|,3)=1\) を仮定しない。

\paragraph{設定.}
\(r\ge1\)、\(m\ge0\)、\(h_i^6\ne-27\) とし、
\[
 \begin{aligned}
 A&=\bigotimes_{i=1}^rS_{h_i}\otimes_\CC T,
 &T&=\CC[u_1,\ldots,u_m],\\
 W&=\sum_{i=1}^r g_i+P,
 &R&=A/(W)
 \end{aligned}
\]
と置く。ここで \(g_i=G_{h_i}\)、\(P\in T_3\) であり、
\(m>0\) の場合は \(T/(\partial_1P,\ldots,\partial_mP)\) が有限次元、
\(m=0\) の場合は \(T=\CC,P=0\) とする。
全ての \(\sigma_i\) は有限位数 \(q_i\) を持つものとする。
\(h_i=0\) は多項式環のfactorで、\(q_i=1\) と解釈する。
周囲の代数 \(A\) はnoetherian AS-regularであり、\(W\) は中心正則元である
（本文のテンソル積の正則性の命題）。
従ってMFと特異圏の比較 \cite[Theorem 5.8]{CCKM} を用いられる。
以下ではこれらの圏を冪等完備化し、\(k=R/R_{>0}\) は特異圏での像、
またはその有限自由stabilizationも表す。

\(s=r\) (\(m=0\))、\(s=r+1\) (\(m>0\)) と置き、
\[
 \Gamma_s=\ZZ^s/\langle3e_i-3e_1:2\le i\le s\rangle,
 \qquad w=3e_i
\]
とする。各 \(S_{h_i}\) の生成元の次数を \(e_i\) とし、
中心変数があれば全ての \(u_a\) の次数を \(e_s\) とする。
通常次数は \(\Gamma_s\to\ZZ,\ e_i\mapsto1\) による粗化であり、このとき \(w=3\) である。
この写像は有限の核を持ち、正の通常次数を与えるため、
\(\Gamma_s\)-斉次自由分解と斉次射を用いる同じMF--特異圏比較が適用できる。

\begin{proposition}[PIの場合のpropernessと残余体生成]
\label{prop:notes-pi-residue-generator}
上の仮定の下で \(\MF^{\Gamma_s}(A,W)\) と \(\MF^\ZZ(A,W)\) はproperである。
特異圏において
\[
 \begin{aligned}
 \Dsg^{\Gamma_s}(R)&=\thick(K_{\Gamma_s}),\\
 K_{\Gamma_s}&=
 \bigoplus_{(a_1,\ldots,a_s)\in\{0,1,2\}^s}
 k(a_1e_1+\cdots+a_se_s)
 \end{aligned}
 \tag{N1}\label{eq:notes-pi-gamma-generator}
\]
および
\[
 \Dsg^\ZZ(R)=\thick(K_{\ZZ}),\qquad
 K_{\ZZ}=k\oplus k(1)\oplus k(2)
 \tag{N2}\label{eq:notes-pi-z-generator}
\]
が成立する。
\end{proposition}

\begin{proof}
\paragraph{一次資料から用いる中心の事実.}
\(Z_i=Z(S_{h_i})\)、\(Y_i=\Spec Z_i\) とし、\(o_i\) を斉次頂点とする。
\(h_i\ne0\) のとき、\(S_{h_i}\) が有限生成可換中心 \(Z_i\) 上有限であることは
\cite[Theorem 3.7(c), p.~38]{SmithTate} による。
さらに \cite[Theorem 1.3(1),(2a)]{WWY3} は
\[
 V(g_i)\setminus\{o_i\}\subset Y_i^{\mathrm{sm}}
       =\operatorname{Az}(S_{h_i}),\qquad
 V(g_i)\setminus\{o_i\}\ \text{は滑らか}
 \tag{N3}\label{eq:notes-pi-zero-slice}
\]
を与える。原論文の三次中心元 \(g\) と \(g_i\) は非零定数倍なので、
零sliceについて同じ主張になる。
従って \(y_i\in V(g_i)\setminus\{o_i\}\) において
\[
 \widehat{(Z_i)_{y_i}}\simeq\CC[[g_i,a,b]],\qquad
 S_{h_i}\otimes_{Z_i}\widehat{(Z_i)_{y_i}}
       \simeq M_{q_i}(\CC[[g_i,a,b]]).
 \tag{N4}\label{eq:notes-pi-matrix-completion}
\]
実際、滑らかな三次元中心と滑らかな零sliceから \(dg_i\ne0\) が従い、
\(g_i\) を正則パラメータに選べる。
完全局所環上のAzumaya代数は、剰余体 \(\CC\) 上でsplitし、
行列冪等元の持ち上げによって行列環になる。
\(h_i=0\) でも、行列sizeを一とし、Fermat三次式の勾配を用いれば
\eqref{eq:notes-pi-zero-slice}--\eqref{eq:notes-pi-matrix-completion} が成立する。

\paragraph{行列Homの支持.}
\(Z=\bigotimes_i Z_i\otimes T\) は有限型可換代数であり、\(A\) は \(Z\) 上有限である。
有限自由factorization \(Q,Q'\) の次数を忘れた二周期行列Hom複体を
\[
 \overline C=\Hom_A^{\ZZ/2}(Q,Q'),\qquad
 \partial(f)=d_{Q'}f-(-1)^{|f|}fd_Q
\]
と書く。その二つの項は有限生成 \(Z\)-加群なので、
\(\overline H^\epsilon=H^{\bar\epsilon}(\overline C)\) も有限生成である。
ブロックEuler導分と中心偏微分により
\[
 d_Q\delta_i(d_Q)+\delta_i(d_Q)d_Q=3g_i\Id_Q,\qquad
 d_Q\partial_{u_a}(d_Q)+\partial_{u_a}(d_Q)d_Q
       =(\partial_aP)\Id_Q.
\]
これらをHomへ合成することで、\(g_i\) と \(\partial_aP\) は
\(\overline H^\epsilon\) を殺す。
Jacobian idealのradicalは \((u_1,\ldots,u_m)\) なので
\[
 \operatorname{Supp}_Z\overline H^\epsilon
 \subset V(g_1,\ldots,g_r,u_1,\ldots,u_m).
 \tag{N5}\label{eq:notes-pi-support}
\]

右辺の閉点 \(y=(y_1,\ldots,y_r,0)\) で、ある \(y_i\ne o_i\) とする。
他のSklyanin factorのテンソル積をその中心点で完備化した代数を \(B\) とすると、
\eqref{eq:notes-pi-matrix-completion} と中心上の完備テンソル積により
\[
 \widehat A_y\simeq
 M_{q_i}\!\left(B[[g_i,a,b,u_1,\ldots,u_m]]\right),
 \qquad W=g_i+\sum_{j\ne i}g_j+P(u).
 \tag{N6}\label{eq:notes-pi-local-matrix}
\]
ここでは次数を忘れてよい。
右辺の形式変数 \(g_i\) に関する成分ごとの導分を \(\partial_{g_i}\) とすれば、
\(\partial_{g_i}W=1\) である。
任意の完備化したfactorizationの微分 \(d\) に対して
\[
 h=\partial_{g_i}(d),\qquad dh+hd=\Id
 \tag{N7}\label{eq:notes-pi-contraction}
\]
となる。従って局所的なungraded MFは実際にcontractibleである。
行列Homにも \(\mathcal H(f)=h_{Q'}f\) による
\(\partial\mathcal H+\mathcal H\partial=\Id\) がある。
中心完備化はflatなのでHom cohomologyと可換であり、
局所環上ではfaithfully flatなので
\((\overline H^\epsilon)_y=0\) が従う。

こうして \(\overline H^\epsilon\) の支持は唯一の頂点
\(o=(o_1,\ldots,o_r,0)\) に含まれる。
\(Z\) は有限型可換 \(\CC\)-代数であるから、
\(\overline H^0\oplus\overline H^1\) は有限次元 \(\CC\)-ベクトル空間である。
この議論には、完備化した他のfactorのglobal dimensionを求める必要はない。

\paragraph{周期性とproperness.}
二周期Homを、偶数部分の微分が内部次数零、奇数部分の微分が内部次数 \(w\)
となるように次数付ける。
具体的に \(Q=(Q_1\to Q_0\to Q_1(w))\) の規約では
\[
 \begin{aligned}
 \overline C^0&=\underline{\Hom}_A(Q_0,Q'_0)
       \oplus\underline{\Hom}_A(Q_1,Q'_1),\\
 \overline C^1&=\underline{\Hom}_A(Q_0,Q'_1(w))
       \oplus\underline{\Hom}_A(Q_1,Q'_0).
 \end{aligned}
\]
とする。ここで \(\underline{\Hom}\) は全ての内部次数のHomの直和である。
\(\overline H^\epsilon\) はこの内部次数を継承し、\(\Lambda=\Gamma_s,\ZZ\) のいずれでも
\[
 \begin{aligned}
 H^{2n}\Hom_{\MF^\Lambda}(Q,Q')
      &=(\overline H^0)_{nw},\\
 H^{2n+1}\Hom_{\MF^\Lambda}(Q,Q')
      &=(\overline H^1)_{nw}.
 \end{aligned}
 \tag{N8}\label{eq:notes-pi-hom-periodicity}
\]
\(w\) は無限位数であり、有限次元の \(\overline H^\epsilon\) は有限個の内部次数しか持たない。
従ってMFのHomの全cohomologyの直和は有限次元である。
有限個のconeと直和因子にもこの結論は保たれるので、冪等完備化したMF圏はproperである。

\paragraph{残余体による生成.}
任意の有限自由factorization \(Q\) を取る。
中心の正次数斉次生成元を \(z_1,\ldots,z_t\) とすると、
全内部次数を含むstable endomorphism加群が頂点に支持されることから、十分大きい \(d\) に対して
\(z_j^d\Id_Q\) は全て零ホモトピックである。
このclassの内部次数は \(d\deg z_j\) であり、
Homの対応する斉次成分を取ることで、
\[
 z_j^d:Q(-d\deg z_j)\longrightarrow Q
\]
のnull-homotopyを次数を保った形に選べる。
中心Koszul複体とのtensor積はこれらのhomotopyを保つので、
これらの乗法のconeを順次取った対象には \(Q\) が直和因子として入る。

\(\Dsg^\Lambda(R)\) で \(Q\) を表す有限生成moduleを \(M\) とする。
今のconeの反復は、boundedな中心Koszul複体
\(K_R(z_1^d,\ldots,z_t^d)\otimes_R M\) の特異圏での像である。
そのcohomologyは \((z_1^d,\ldots,z_t^d)\) に殺される。
\(R\) は \(Z\) 上有限で、\(Z/(z_1^d,\ldots,z_t^d)\) は有限次元なので、
これらのcohomologyも有限次元である。
有限次元の次数付き \(R\)-moduleは正の通常次数によるfiltrationを持ち、
各商は有限個の \(k(\lambda)\) の直和である。
従って全ての \(Q\) は \(\{k(\lambda)\}_{\lambda\in\Lambda}\) で生成される。
冪等完備化についても同じ結論になる。

最後に \((w)=[2]\) であり、
\[
 \Gamma_s/\ZZ w\simeq(\ZZ/3)^s,\qquad \ZZ/3\ZZ\simeq\ZZ/3
\]
なので、それぞれの剰余類の代表だけで十分である。
これが \eqref{eq:notes-pi-gamma-generator} と \eqref{eq:notes-pi-z-generator} を与える。
\end{proof}

\paragraph{次の計算への使い方.}
細分次数では \(3^s\) 個、通常次数では \(k,k(1),k(2)\) の三つのstabilizationを用意すればよい。
その直和 \(K_\Lambda\) に対して
\[
 B_\Lambda=\REnd_{\Dsg^\Lambda(R)}(K_\Lambda),\qquad
 \MF^\Lambda(A,W)\simeq_{\mathrm{Morita}}\Perf(B_\Lambda)
\]
となる。自己Homはstabilization後のMFのdg模型で取り、\(H^*(B_\Lambda)\) は有限次元である。
本文の有限位数多ブロック比較と組み合わせると、この模型はsmooth properでもある。
この生成子の記述は、外部積による全生成性の証明とは独立に、
全MF圏を具体的な有限個の対象の自己Homから調べる入口を与える。
既知の中心幾何からの上の帰結について、先行研究上の新規性は未確認である。

